\documentclass{article}
\usepackage{amsmath}
\usepackage{amssymb}
\usepackage{amsthm}
\usepackage{mathtools}
\usepackage{geometry}
\usepackage{hyperref}

\newcommand{\rk}{\operatorname{rank}}
\newcommand{\GL}{\operatorname{GL}}
\newcommand{\R}{\mathbb{R}}
\newcommand{\Mat}{\mathrm{Mat}}
\newcommand{\im}{\operatorname{im}}
\newcommand{\Ker}{\operatorname{Ker}}

\theoremstyle{definition}
\newtheorem{theorem}{Theorem}[section]
\newtheorem{proposition}[theorem]{Proposition}
\newtheorem{lemma}[theorem]{Lemma}
\newtheorem{corollary}[theorem]{Corollary}
\newtheorem{definition}[theorem]{Definition}
\newtheorem{remark}[theorem]{Remark}

\title{Task 3: Invariants of the End-to-End Map}
\date{}

\begin{document}
\maketitle

\section{Goal}

Extract all constraints on the end-to-end map
\[
M=\mu(W)
\]
that are forced by the cyclic equations
\[
W_{<\ell}W_0W_{>\ell}=0.
\]

\section{Questions to answer}

This file should address the following:

\begin{itemize}
\item Which image and kernel constraints on $M$ are visible from the cyclic equations?
\item Which constraints depend only on $M$, and which still depend on the full tuple $W$?
\item Which of these constraints are invariant under the endpoint group $K_S$?
\end{itemize}

\section{Visible constraints on $M$}

\begin{proposition}[Endpoint annihilation identities]\label{prop:endpoint-annihilation}
Let $W\in \Gamma_{d,S}$ and set
\[
M:=\mu(W).
\]
Then
\[
W_0M=0
\qquad\text{and}\qquad
MW_0=0.
\]
\end{proposition}

\begin{proof}
Since $W\in \Gamma_{d,S}$, the cyclic equations
\[
E_\ell(W):=W_{<\ell}W_0W_{>\ell}=0
\]
hold for every $1\le \ell\le L$.

For $\ell=1$ one has
\[
E_1(W)=W_0W_{>1}=0.
\]
Multiplying on the right by $W_1$ gives
\[
W_0W_{>1}W_1=W_0M=0.
\]

For $\ell=L$ one has
\[
E_L(W)=W_{<L}W_0=0.
\]
Multiplying on the left by $W_L$ gives
\[
W_LW_{<L}W_0=MW_0=0.
\]
\end{proof}

\begin{corollary}[Image and kernel constraints]\label{cor:image-kernel-constraints}
If $W\in \Gamma_{d,S}$ and $M=\mu(W)$, then
\[
\im(M)\subseteq \Ker(W_0)
\qquad\text{and}\qquad
\im(W_0)\subseteq \Ker(M).
\]
\end{corollary}

\begin{proof}
The identity $W_0M=0$ means exactly that every vector in $\im(M)$ is killed by $W_0$, hence
\[
\im(M)\subseteq \Ker(W_0).
\]
Similarly, $MW_0=0$ means that every vector in $\im(W_0)$ is killed by $M$, hence
\[
\im(W_0)\subseteq \Ker(M).
\]
\end{proof}

\begin{corollary}[Rank-$r$ sharpening under the full-rank data assumption]\label{cor:rank-r-sharpening}
Assume that $\Sigma_{XY}$ has full rank $d_L$ and let $W\in \Gamma_{d,S}^r$ with $M=\mu(W)$. Then
\[
\Ker(W_0)=\im(P_S),
\qquad
\rk(W_0)=d_L-r,
\]
and therefore
\[
\im(M)=\Ker(W_0)=\im(P_S).
\]
\end{corollary}

\begin{proof}
Recall that
\[
W_0=\Sigma_{XY}P_Q.
\]
Because $\Sigma_{XY}$ has full rank $d_L$, the linear map $\Sigma_{XY}:\R^{d_L}\to \R^{d_0}$ is injective. Therefore
\[
\Ker(W_0)=\Ker(P_Q)=\im(P_S),
\]
and
\[
\rk(W_0)=\rk(P_Q)=d_L-r.
\]

By Corollary~\ref{cor:image-kernel-constraints},
\[
\im(M)\subseteq \Ker(W_0)=\im(P_S).
\]
But $W\in \Gamma_{d,S}^r$ implies $\rk(M)=r$, so $\dim \im(M)=r$. Since $\im(P_S)$ also has dimension $r$, the inclusion must be an equality:
\[
\im(M)=\im(P_S).
\]
\end{proof}

\begin{proposition}[The visible constraints are $K_S$-invariant]\label{prop:ks-invariance}
Let $(g_0,g_L)\in K_S$ and let
\[
M':=g_LMg_0^{-1}.
\]
Then:
\begin{enumerate}
\item $\rk(M')=\rk(M)$;
\item if $W_0M=0$, then $W_0M'=0$;
\item if $MW_0=0$, then $M'W_0=0$.
\end{enumerate}
\end{proposition}

\begin{proof}
Since $g_0$ and $g_L$ are invertible,
\[
\rk(M')=\rk(g_LMg_0^{-1})=\rk(M).
\]

Next, because $(g_0,g_L)\in K_S$, one has
\[
g_0W_0g_L^{-1}=W_0.
\]
Equivalently,
\[
W_0g_L=g_0W_0
\qquad\text{and}\qquad
g_0^{-1}W_0=W_0g_L^{-1}.
\]
Hence
\[
W_0M'=W_0g_LMg_0^{-1}=g_0W_0Mg_0^{-1}.
\]
So if $W_0M=0$, then $W_0M'=0$.

Similarly,
\[
M'W_0=g_LMg_0^{-1}W_0=g_LMW_0g_L^{-1}.
\]
So if $MW_0=0$, then $M'W_0=0$.
\end{proof}

\section{What does not descend to $M$ alone}

\begin{proposition}[The remaining cyclic information lives in the factorization]\label{prop:factorized-information}
Membership in the cyclic locus is not determined by the end-to-end map alone. More precisely, there exist two tuples
\[
W,\widetilde W\in \Omega
\]
with the same end-to-end map
\[
\mu(W)=\mu(\widetilde W),
\]
such that
\[
W\in \Gamma_{d,S}^r
\qquad\text{but}\qquad
\widetilde W\notin \Gamma_{d,S}^r.
\]
\end{proposition}

\begin{proof}
Take
\[
L=3,
\qquad
d_0=d_1=d_2=d_3=2,
\qquad
\Sigma_{XY}=I_2,
\qquad
P_Q=
\begin{pmatrix}
0&0\\
0&1
\end{pmatrix}.
\]
Then
\[
W_0=\Sigma_{XY}P_Q=
\begin{pmatrix}
0&0\\
0&1
\end{pmatrix},
\qquad
P_S=
\begin{pmatrix}
1&0\\
0&0
\end{pmatrix},
\qquad
r=1.
\]
Set
\[
M:=
\begin{pmatrix}
1&0\\
0&0
\end{pmatrix}.
\]
Then
\[
\rk(M)=1,
\qquad
W_0M=0,
\qquad
MW_0=0.
\]

Now define
\[
W_1=M,\qquad W_2=M,\qquad W_3=I_2.
\]
Since $M^2=M$, we have
\[
\mu(W)=W_3W_2W_1=I_2MM=M.
\]
Moreover,
\begin{align*}
E_1(W)&=W_0W_3W_2=W_0M=0, \\
E_2(W)&=W_1W_0W_3=MW_0=0, \\
E_3(W)&=W_2W_1W_0=MMW_0=MW_0=0.
\end{align*}
So $W\in \Gamma_{d,S}^1$.

Next define
\[
\widetilde W_1=I_2,\qquad \widetilde W_2=I_2,\qquad \widetilde W_3=M.
\]
Then
\[
\mu(\widetilde W)=\widetilde W_3\widetilde W_2\widetilde W_1=M.
\]
But
\[
E_3(\widetilde W)=\widetilde W_2\widetilde W_1W_0=I_2I_2W_0=W_0\neq 0.
\]
Hence
\[
\widetilde W\notin \Gamma_{d,S}^1.
\]
Thus the same matrix $M$ admits both a cyclic and a non-cyclic factorization.
\end{proof}

\begin{corollary}[Separation lemma for Task 3]\label{cor:separation-lemma}
The cyclic equations impose two kinds of information:
\begin{enumerate}
\item product-level constraints visible on $M=\mu(W)$, namely at least
\[
\rk(M)=r,
\qquad
W_0M=0,
\qquad
MW_0=0;
\]
\item additional factorization-level constraints, encoded by the intermediate equations
\[
W_{<\ell}W_0W_{>\ell}=0
\qquad (1<\ell<L),
\]
which are not determined by $M$ alone.
\end{enumerate}
\end{corollary}

\begin{remark}
Task~3 isolates the universal constraints on $M$ that are visibly forced by the cyclic equations. It does \emph{not} yet prove that these constraints characterize the relevant $K_S$-orbit of $P_SW^\ast$. That is the remaining burden of Task~4.
\end{remark}

\section{Deliverables}

This task is complete: the product-level constraints forced by the cyclic equations are identified, their $K_S$-invariance is proved, and an explicit rank-one example shows that the remaining cyclic information still lives in the factorization.

\end{document}
